\chapter*{Preface}
\addcontentsline{toc}{chapter}{Preface}

There is a particular kind of frustration that brings people to cryptography. You read that TLS ``encrypts your connection,'' that Signal offers ``end-to-end encryption,'' that your WiFi password is ``protected by WPA3,'' and somewhere underneath all of it sits mathematics you were told to trust but never got to see. The padlock in the browser is a promise. This book is about what backs the promise.

I wrote it because the books that taught me cryptography each held one corner of the elephant. Some gave me clean theorems and security proofs but never showed a packet on a wire. Some gave me practical recipes -- use AES-GCM, never roll your own crypto -- without ever explaining \emph{why} AES-GCM and not something else, or what ``rolling your own'' actually breaks. And some gave me beautiful intuition with hand-waving exactly where I most wanted rigor. The goal here is to hold all three corners at once: the math, the intuition, and the real-world use, side by side on the same page.

\section*{What this book is, and who it is for}

You are smart, motivated, and willing to do arithmetic by hand when it teaches you something. You do not need to already know number theory or abstract algebra -- we build those from the ground up, which is the guiding philosophy of this book. You should be comfortable with high-school algebra, with the idea of a function, and with the patience to follow a small calculation to its end. If you have written a few hundred lines of code in any language, you will get more out of the toy implementations, but they are not strictly required.

By the end you should be able to start work as an entry-level cryptographer, or a security engineer who actually understands the primitives. Concretely, that means you can:

\begin{itemize}
  \item read a protocol specification (a TLS RFC, the Signal documentation, the WPA3 standard) and follow the cryptographic reasoning without it being a wall of acronyms;
  \item look at a scheme and state its threat model -- who the attacker is, what they can see, what they can change, and what ``secure'' is supposed to mean;
  \item do the math: modular exponentiation, a tiny RSA key, a point addition on an elliptic curve over a small field, a birthday-bound estimate, all by hand on small numbers;
  \item recognize the classic failure modes -- nonce reuse, padding oracles, timing leaks, missing authentication -- and explain the attack, not just name it;
  \item choose the right tool (authenticated encryption, not a raw block cipher; a key-derivation function, not a bare hash) and say why.
\end{itemize}

That last skill is the real bar. A cryptographer is not someone who has memorized algorithms; it is someone who can reason about what an adversary can do, and prove -- or at least argue carefully -- that the adversary cannot do the thing you care about. Everything in this book is bent toward that habit of mind.

\section*{How to read this book}

Each scheme is introduced the same way, and once you see the pattern you can use it as scaffolding for your own understanding.

First comes the \emph{story}. Alice wants to send Bob a message; Eve is listening on the wire, and sometimes Mallory is actively tampering with it. We start there because every cryptographic protocol is a contest between parties with conflicting goals, and the algorithms are the moves in that contest. Naming the players keeps you precise about who knows what.

\begin{marginnotebox}
The cast you will meet over and over: \textbf{Alice} and \textbf{Bob} want to communicate. \textbf{Eve} is a passive eavesdropper -- she reads but does not change. \textbf{Mallory} is an active, in-the-middle attacker -- she can drop, replay, reorder, and forge messages. \textbf{Trent} is a trusted third party (a certificate authority, say). When a scheme protects against Eve but not Mallory, that distinction is usually where the interesting attacks live.
\end{marginnotebox}

Then come \emph{tiny numbers}. Before any general formula, we trace the idea through an example small enough to check on paper. RSA with primes $p = 61$ and $q = 53$. A four-bit toy block cipher you can encrypt by hand. An elliptic curve over $\F_{11}$ with so few points you can list them all. Real little numbers are not a teaching gimmick -- watching $7^{13} \bmod 23$ collapse to a single value is how the abstraction stops being abstract.

Then the \emph{precise math}, and a \emph{security argument}. We state what the scheme is, what assumption it rests on (factoring is hard, the discrete logarithm is hard, this hash behaves like a random function), and we argue why breaking the scheme would mean breaking the assumption. Some of these arguments are full reductions in the style of Katz and Lindell; others are the informal-but-correct sketches you would give a colleague. Either way you should finish a section able to answer ``why do we believe this is secure?'' with more than ``because the textbook said so.''

You do not have to read strictly front to back, but the early chapters genuinely are foundations. The arithmetic and algebra of Part~I are the vocabulary every later chapter speaks. If you already know modular arithmetic and groups cold, skim them and start at the symmetric ciphers; if you do not, resist the urge to skip -- you will only end up back there confused.

\subsection*{The colored boxes, and what each one signals}

The text uses recurring colored boxes so you can tell at a glance what kind of thinking a passage demands. They are part of the book's grammar:

\begin{intuitionbox}
The mental picture -- the analogy or the one-sentence ``aha'' that makes the formal statement feel inevitable. If you remember nothing else from a section, remember the intuition boxes.
\end{intuitionbox}

\begin{whybox}
The argument for \emph{why} something is secure, or why an attack succeeds. This is where the security reasoning lives -- reductions, indistinguishability arguments, the logic that connects an assumption to a guarantee.
\end{whybox}

\begin{formalbox}
The precise definition or theorem, stated the way you would write it in a paper. When intuition and rigor diverge, this box is the authority.
\end{formalbox}

\begin{attackbox}
A concrete break. Not ``this could be weak'' but ``here is exactly how Mallory wins, step by step.'' Attacks teach defense better than any list of best practices, because you cannot defend a boundary you have never seen crossed.
\end{attackbox}

\begin{focusbox}
Key ideas and guiding themes. A compass heading for the chapter so you never lose the forest for the trees.
\end{focusbox}

\begin{implementbox}
Short toy code -- usually Python -- that turns an idea into something you can run. \textbf{Read this warning once and carry it through the whole book:} the code in these boxes is written to be \emph{transparent}, not safe. It leaks timing, it uses weak randomness, it skips the checks that production code must never skip. It exists so you can see the mechanism. \emph{Never} paste it into anything real. Real cryptographic code belongs to vetted, constant-time, well-maintained libraries, and one of the things this book teaches is the judgment to know why.
\end{implementbox}

\begin{applicationbox}
Where the idea shows up in systems you use every day -- TLS, WiFi, Signal, SSH, the cellular network. This is the answer to ``so what?'' and it is also where abstract guarantees meet messy engineering reality.
\end{applicationbox}

\begin{connectionbox}
A pointer to the companion ground-up readers on probability and on algorithms. When we need a birthday bound, the cost of modular exponentiation, or the analysis of a randomized algorithm, the deep version lives there and we link to it rather than reprove it.
\end{connectionbox}

You will also meet \texttt{historybox} (historical context and lessons from past failures), \texttt{walkbox} (a walking question to chew on away from the page), \texttt{trapbox} (a mistake almost everyone makes once), and \texttt{recoverybox} (a place to pause and consolidate). The boxes are signposts, not interruptions -- let them tell you which gear to be in.

\section*{The security mindset}

Before any specific scheme, you need a way of thinking, and it is the hardest thing to teach because it runs against ordinary engineering instincts. An ordinary engineer asks ``does it work?'' -- does the system do the right thing when used the right way. A cryptographer asks the adversary's question: ``how does it fail when someone is \emph{trying} to make it fail?'' The user is not your model of reality; the attacker is.

This flips the burden of proof. In normal software, a feature is correct until a bug is found. In cryptography, a scheme is \emph{broken} until there is a convincing argument that it is not -- and ``I couldn't break it'' is not such an argument, because the person who breaks it may be cleverer, better funded, or simply have more time than you. That is why the field leans so heavily on definitions and proofs. A security definition is a promise made precise: it says exactly which attacker, with exactly which powers, fails to achieve exactly which goal. Without that precision, ``secure'' is a marketing word.

\begin{intuitionbox}
Think of it as a game. We hand the attacker a generous set of powers -- let her choose messages to be encrypted, let her see ciphertexts, sometimes even let her ask for some decryptions -- and we define a single thing she must not be able to do (tell two ciphertexts apart, say, or forge one valid tag). The scheme is secure if \emph{even with all those powers} her chance of winning is no better than blind guessing, up to a negligible sliver. We will make ``negligible'' precise later; for now it means ``so small that no realistic amount of computation closes the gap.''
\end{intuitionbox}

Two principles follow from this mindset and recur on nearly every page.

\subsection*{Kerckhoffs's principle: the secret is the key, never the design}

In 1883 the Dutch linguist Auguste Kerckhoffs laid down a rule for military ciphers that the entire modern field has adopted: a cryptosystem should remain secure even if everything about it is public knowledge \emph{except the key}. The algorithm, the protocol, the source code, the exact wiring -- assume the enemy has all of it. The only thing you get to keep secret is a short string, the key, which is easy to generate fresh, easy to change, and easy to keep in small quantity.

\begin{historybox}
Claude Shannon, founding the mathematical theory of secrecy in 1949, restated Kerckhoffs's 1883 rule so bluntly it became a slogan: ``the enemy knows the system.'' Design as if your adversary has read your code, bought your hardware, and reverse-engineered your firmware -- because eventually one of them will. Historically, governments and vendors resisted this, relying on secret designs known as ``security through obscurity.'' That stance has a long record of failure: the A5/1 cipher protecting GSM phone calls was kept secret for years, but once its design leaked and was analyzed in the open, it fell quickly (as we will dissect in Chapter~29). Open standards like AES earn trust only after thousands of independent cryptanalysts try to break them; today, the hardest frontier is defending open algorithms against hardware side channels like timing and power leaks.
\end{historybox}

\begin{whybox}
Why insist on this? Because secrets that are hard to change are secrets that will eventually leak and then cannot be repaired. An algorithm baked into a billion devices cannot be swapped overnight; a leaked key can be rotated in seconds. Kerckhoffs's principle pushes all the fragility into the one component that is cheap to replace. It also has a social benefit that turns out to be decisive: a public algorithm can be attacked by the whole world's cryptanalysts for years before you trust it. Every cipher in this book that you should actually use -- AES, the SHA-2 and SHA-3 families, Curve25519 -- earned that trust by surviving open, published assault. A secret design earns nothing, because no one outside could ever try to break it.
\end{whybox}

\begin{trapbox}
``We use a proprietary, undisclosed encryption algorithm for extra security.'' Whenever you read this in a product's marketing, translate it as ``no qualified outsider has ever checked our math.'' Secrecy of the design is not a second layer of protection on top of a strong cipher; historically it is a substitute for one, and a poor substitute at that. Extra obscurity is fine as a thin outer wrapper -- it raises the cost of a casual attacker -- but it must never be the thing your security rests on.
\end{trapbox}

\subsection*{Don't roll your own crypto -- and why this book seems to do exactly that}

The second principle sounds like it contradicts the existence of this book: do not write your own cryptographic code for real use. We will write plenty of cryptographic code. The resolution is in the word \emph{real}. Understanding a scheme well enough to implement a toy version is exactly what makes you competent to \emph{use} the real one correctly -- to pick the right mode, supply a proper nonce, check a tag before trusting a plaintext. But the gap between a textbook-correct implementation and a deployment-safe one is enormous, and it is paved with side channels: timing differences, power draw, cache behavior, error messages that leak one bit at a time.

\begin{attackbox}
A textbook RSA decryption can be perfectly correct and still hand an attacker your private key, because the time it takes leaks information about the secret exponent. A textbook AES in CBC mode can be perfectly correct and still surrender the plaintext to a \emph{padding oracle} -- an attacker who learns only whether decryption ``padded correctly'' can peel off the message one byte at a time, without ever recovering the key. We will build these very attacks by hand in later chapters. They are the reason production crypto is constant-time, authenticated, and written by specialists. Seeing them is what immunizes you against shipping them.
\end{attackbox}

So: implement everything here to learn it. Trust none of it in production. Hold both at once. That tension -- deep understanding paired with deep humility about deployment -- is the mark of a real cryptographer, and it is the attitude every \texttt{implementbox} in this book is training.

\section*{The arc: from modular arithmetic to TLS, WiFi, and Signal}

The book climbs a single long staircase. Each step rests on the ones below it, and it helps to see the whole flight before you start.

\textbf{Part~I -- Foundations.} We begin with the mathematics: modular arithmetic, primes and the greatest common divisor, groups, rings, fields, and the finite fields that the rest of the subject is built on. Then probability and a first taste of information theory, and the idea of computational hardness -- the notion that some problems are easy to state and effectively impossible to solve at scale, which is the bedrock under all of public-key cryptography. This part has the least to do with secrets and the most to do with whether anything later will make sense. Take it slowly.

\textbf{Part~II -- Symmetric cryptography.} Alice and Bob already share a secret key; what can they do with it? We reach perfect secrecy and the one-time pad -- the one scheme that is provably unbreakable, and a lovely lesson in why provably unbreakable is not the same as practical. Then stream ciphers, block ciphers (the Feistel construction, DES, and AES), the modes that turn a block cipher into something that can encrypt a whole message, hash functions, message authentication codes and HMAC, and finally authenticated encryption -- the modern default that gives you confidentiality and integrity together, which it turns out you almost always need at the same time.

\textbf{Part~III -- Public-key cryptography.} Now we drop the assumption that Alice and Bob already share a key, and confront the question that launched modern cryptography: how do two strangers agree on a secret over a wire that an enemy is reading? Diffie--Hellman key exchange, RSA, the discrete logarithm and ElGamal, elliptic curves -- the same hardness assumptions from Part~I, now turned into working machinery. Digital signatures, key exchange with forward secrecy, and the certificate-and-PKI scaffolding that lets your browser decide whom to trust.

\textbf{Part~IV -- Real systems.} Here the staircase reaches the buildings you actually live in. TLS, the protocol behind every \texttt{https}. The evolution of WiFi security up through WPA3. Signal's double ratchet, which gives messaging both forward secrecy and post-compromise recovery. SSH, the cryptography woven through the cellular network, VPNs and tunnels. These chapters are where every primitive from the earlier parts gets composed into a system -- and where you learn that composing secure pieces does not automatically yield a secure whole.

\textbf{Part~V -- The frontier and the craft.} Zero-knowledge proofs, commitments and secret sharing, homomorphic encryption, multiparty computation, post-quantum cryptography, and the cryptography under blockchains. Then the practitioner's chapters that tie the whole book together: side-channel and timing attacks, the catalog of common pitfalls, and how to implement cryptography safely in the real world. Two appendices collect the number theory and algebra for reference.

\begin{connectionbox}
This is one of three companion ground-up readers. Where we need the deep version of a probabilistic argument -- the precise birthday bound behind hash collisions, the analysis of a randomized primality test -- the probability book carries the full treatment. Where we need the cost of an algorithm -- how many multiplications modular exponentiation really takes, why fast multiplication matters for big keys -- the algorithms book has it. Cryptography sits exactly at the meeting point of ``what is hard to compute'' and ``what is likely by chance,'' so you will see those pointers often.
\end{connectionbox}

\begin{focusbox}
A single thread runs the length of the staircase: the move from \emph{information-theoretic} security (the one-time pad, secure even against an adversary with infinite computing power, but crippled by the need for a key as long as the message) to \emph{computational} security (everything else, secure only against adversaries with realistic resources, but practical enough to run the internet). Almost all of modern cryptography lives in that second world, and almost all of its subtlety comes from making ``realistic resources'' precise. Keep one eye on that distinction the whole way up.
\end{focusbox}

The dependency map below sketches how the parts lean on one another. You can take many paths through it, but no path skips the foundations.

\begin{center}
\begin{tikzpicture}[
    font=\small,
    node distance=1.1cm and 1.2cm,
    partbox/.style={draw=#1!75!black, fill=#1!10, rounded corners=2pt, align=center, line width=0.7pt, minimum width=3.8cm, minimum height=1.25cm, inner sep=4pt},
    corridor/.style={->, >=stealth, thick, draw=gray!70!black}
  ]
  % Foundations header and box
  \node[partbox=sourceblue] (found) {
    \textbf{\color{sourceblue!90!black}Part~I: Foundations}\\
    \scriptsize Modular arithmetic, groups, fields, hardness
  };

  % Middle tier: Symmetric (left) and Public-Key (right)
  \node[partbox=targetgreen, below left=1.1cm and -1.0cm of found] (sym) {
    \textbf{\color{targetgreen!90!black}Part~II: Symmetric Crypto}\\
    \scriptsize Stream/block ciphers, hashes, MACs, AEAD
  };

  \node[partbox=quotientorange, below right=1.1cm and -1.0cm of found] (pub) {
    \textbf{\color{quotientorange!90!black}Part~III: Public-Key Crypto}\\
    \scriptsize Diffie--Hellman, RSA, ECC, signatures
  };

  % Systems tier
  \node[partbox=structurepurple, below=2.7cm of found] (sys) {
    \textbf{\color{structurepurple!90!black}Part~IV: Real Protocols}\\
    \scriptsize TLS 1.3, Wi-Fi (WPA3), Signal ratchet, SSH
  };

  % Advanced / Practice tier
  \node[partbox=cyan, below=0.9cm of sys] (front) {
    \textbf{\color{cyan!80!black}Part~V: Advanced \& Safe Implementation}\\
    \scriptsize Zero-knowledge, post-quantum, timing leaks \& defenses
  };

  % Dedicated arrow corridors with generous clearance
  \draw[corridor] (found.south) -- ++(0,-0.3) -| (sym.north);
  \draw[corridor] (found.south) -- ++(0,-0.3) -| (pub.north);
  \draw[corridor] (sym.south) |- (sys.west);
  \draw[corridor] (pub.south) |- (sys.east);
  \draw[corridor] (sys.south) -- (front.north);
\end{tikzpicture}
\end{center}

\section*{One worked example, so you know what ``by hand'' means}

To set expectations, here is the flavor of the small-number work you will do throughout, using a piece of arithmetic that powers nearly every public-key scheme: modular exponentiation, computed the smart way.

Suppose we want $7^{13} \bmod 23$. The naive route -- multiply seven by itself thirteen times, then reduce -- works but scales terribly; with the thousand-bit exponents of real cryptography it would never finish. The trick is \emph{square-and-multiply}. Write the exponent in binary: $13 = 1101_2 = 8 + 4 + 1$. Then
\[
7^{13} = 7^{8} \cdot 7^{4} \cdot 7^{1},
\]
and we get the high powers by repeated squaring, reducing mod $23$ at every step so the numbers never grow:
\[
7^{1} \equiv 7,\qquad
7^{2} \equiv 49 \equiv 3,\qquad
7^{4} \equiv 3^{2} = 9,\qquad
7^{8} \equiv 9^{2} = 81 \equiv 81 - 69 = 12 \pmod{23}.
\]
Now multiply together the pieces the binary digits selected (the $8$, the $4$, and the $1$):
\[
7^{13} \equiv 7^{8}\cdot 7^{4}\cdot 7^{1} \equiv 12 \cdot 9 \cdot 7 \pmod{23}.
\]
Take it in steps: $12 \cdot 9 = 108 \equiv 108 - 92 = 16$, and $16 \cdot 7 = 112 \equiv 112 - 92 = 20$. So
\[
7^{13} \equiv 20 \pmod{23}.
\]

\begin{center}
\begin{tikzpicture}[
    font=\small,
    >=stealth,
    lane/.style={draw=sourceblue!70, fill=sourceblue!8, rounded corners=2pt, minimum height=0.7cm, align=center, inner sep=3pt},
    gate/.style={draw=quotientorange!80, fill=quotientorange!12, rounded corners=2pt, minimum height=0.7cm, align=center, inner sep=3pt},
    ruler/.style={font=\footnotesize\ttfamily, text=gray!70!black, anchor=east}
  ]
  % Upper lane: Bit powers
  \node[ruler] at (-1.5, 1.3) {bit index:};
  \node[font=\footnotesize\ttfamily, text=gray!70!black] at (0, 1.3) {bit 3 ($2^3=8$)};
  \node[font=\footnotesize\ttfamily, text=gray!70!black] at (3.2, 1.3) {bit 2 ($2^2=4$)};
  \node[font=\footnotesize\ttfamily, text=gray!70!black] at (6.4, 1.3) {bit 1 ($2^1=2$)};
  \node[font=\footnotesize\ttfamily, text=gray!70!black] at (9.6, 1.3) {bit 0 ($2^0=1$)};

  % Exponent bit cells
  \node[ruler] at (-1.5, 0.45) {binary $13$:};
  \node[lane, minimum width=2.3cm] (b3) at (0, 0.45) {\textbf{1} $\to$ fold in $7^8$};
  \node[lane, minimum width=2.3cm] (b2) at (3.2, 0.45) {\textbf{1} $\to$ fold in $7^4$};
  \node[lane, fill=gray!8, draw=gray!40, text=gray!60!black, minimum width=2.3cm] (b1) at (6.4, 0.45) {\textbf{0} $\to$ skip};
  \node[lane, minimum width=2.3cm] (b0) at (9.6, 0.45) {\textbf{1} $\to$ fold in $7^1$};

  % Lower lane: Repeated squarings
  \node[ruler] at (-1.5, -0.65) {squaring mod 23:};
  \node[gate, minimum width=2.3cm] (p3) at (0, -0.65) {$7^8 \equiv 12$};
  \node[gate, minimum width=2.3cm] (p2) at (3.2, -0.65) {$7^4 \equiv 9$};
  \node[gate, minimum width=2.3cm] (p1) at (6.4, -0.65) {$7^2 \equiv 3$};
  \node[gate, minimum width=2.3cm] (p0) at (9.6, -0.65) {$7^1 \equiv 7$};

  % Squaring arrows right-to-left
  \draw[->, thick, quotientorange!80] (p0.west) -- node[above=2pt, font=\scriptsize] {square} (p1.east);
  \draw[->, thick, quotientorange!80] (p1.west) -- node[above=2pt, font=\scriptsize] {square} (p2.east);
  \draw[->, thick, quotientorange!80] (p2.west) -- node[above=2pt, font=\scriptsize] {square} (p3.east);

  % Bottom summary
  \node[draw=targetgreen!75!black, fill=targetgreen!10, rounded corners=2pt, font=\small, minimum width=11.5cm, minimum height=0.7cm] at (4.8, -1.7) {
    Accumulator product: $(7^8) \cdot (7^4) \cdot (7^1) \equiv 12 \cdot 9 \cdot 7 \equiv 16 \cdot 7 \equiv \mathbf{20} \pmod{23}$
  };
\end{tikzpicture}
\end{center}

Four squarings and two multiplications replace twelve sequential multiplies. That logarithmic growth makes 2048-bit modular exponentiations run in a matter of milliseconds.

\begin{implementbox}
The same algorithm, written to be read rather than deployed. Python already has a fast, safe \texttt{pow(base, exp, mod)} built in -- this version exists only to make the squaring visible.

\begin{lstlisting}[language=Python]
def modexp(base, exp, mod):
    # Square-and-multiply. LEARNING ONLY:
    # not constant-time -- leaks the exponent's bits via timing.
    result = 1
    base = base % mod
    while exp > 0:
        if exp & 1:                  # low bit is 1 -> fold in this power
            result = (result * base) % mod
        base = (base * base) % mod    # square for the next bit
        exp >>= 1                     # move to the next bit of the exponent
    return result

print(modexp(7, 13, 23))   # -> 20
\end{lstlisting}

The comment is not a formality. The branch \texttt{if exp \& 1} runs a multiplication only when the exponent bit is set, so the running time depends on the secret exponent -- exactly the leak that lets an attacker recover a private key by timing. Real libraries spend serious effort erasing that dependence. This is your first encounter with the gap between ``correct'' and ``safe,'' and it will not be your last.
\end{implementbox}

\begin{connectionbox}
Why is square-and-multiply the right algorithm and not merely a clever one? Because it runs in time proportional to the number of \emph{bits} in the exponent, not its magnitude -- logarithmic in the exponent itself. The companion algorithms reader develops this cost model carefully and explains why that logarithmic gap is the difference between cryptography that runs in milliseconds and cryptography that never finishes.
\end{connectionbox}

That single operation -- raise to a power, reduce mod $n$ -- is, astonishingly, almost the entire engine of RSA. To make the point concrete, and to show you that ``by hand'' really does mean by hand, here is a complete RSA encryption and decryption on numbers small enough to check on a napkin. Take the two primes the running RSA example will use, $p = 61$ and $q = 53$. Then
\[
n = pq = 61 \cdot 53 = 3233, \qquad
\varphi(n) = (p-1)(q-1) = 60 \cdot 52 = 3120,
\]
where $\varphi(n)$ counts the integers below $n$ that share no factor with it -- the size of the group we will be exponentiating in. Pick a public exponent $e = 17$, which is coprime to $3120$. The private exponent $d$ is the modular inverse of $e$ modulo $\varphi(n)$, the number with $e\,d \equiv 1 \pmod{3120}$; the extended Euclidean algorithm (which we develop in full in Part~I) returns $d = 2753$. You can check the defining relation directly: $17 \cdot 2753 = 46801 = 15 \cdot 3120 + 1$, so $17 \cdot 2753 \equiv 1 \pmod{3120}$ as required.

Now encrypt the message $m = 65$ (think of it as a single byte). The ciphertext is
\[
c \equiv m^{e} \equiv 65^{17} \pmod{3233},
\]
which is exactly a square-and-multiply away. Writing $17 = 10001_2 = 16 + 1$ and squaring mod $3233$:
\[
65^{1} \equiv 65,\quad
65^{2} \equiv 4225 \equiv 992,\quad
65^{4} \equiv 992^{2} = 984064 \equiv 1232,\quad
65^{8} \equiv 1232^{2} = 1517824 \equiv 1547,
\]
and one more squaring gives $65^{16} \equiv 1547^{2} = 2393209 \equiv 789 \pmod{3233}$. The exponent $17$ selects the $16$ and the $1$, so
\[
c \equiv 65^{16}\cdot 65^{1} \equiv 789 \cdot 65 = 51285 \equiv 2790 \pmod{3233}.
\]
To decrypt, raise the ciphertext to the private exponent: $m \equiv c^{d} \equiv 2790^{2753} \pmod{3233}$. By the same square-and-multiply (now over twelve bits) this returns $2790^{2753} \equiv 65$, recovering the original byte. The round trip below puts the whole toy RSA cycle on one line: the public key $(n, e)$ locks the message, and only the modular inverse $d \equiv e^{-1} \pmod{\varphi(n)}$ unlocks it.

\begin{center}
\begin{tikzpicture}[
    font=\small,
    >=stealth,
    box/.style={draw=#1!80!black, fill=#1!10, rounded corners=3pt, thick, align=center, inner sep=5pt, minimum height=1.1cm},
    keytag/.style={draw=#1!80!black, fill=#1!15, rounded corners=2pt, font=\scriptsize, align=center, inner sep=3.5pt}
  ]
  \node[box=sourceblue, minimum width=2.2cm] (pt) at (0, 0) {Plaintext\\$m = 65$};
  \node[box=quotientorange, minimum width=2.5cm] (ct) at (5.5, 0) {Ciphertext\\$c = 2790$};
  \node[box=targetgreen, minimum width=2.2cm] (rec) at (11.0, 0) {Recovered\\$m = 65$};

  \draw[->, thick, sourceblue!85!black] (pt.east) --
    node[above=3pt, font=\scriptsize\bfseries] {$65^{17} \bmod 3233$}
    (ct.west);

  \draw[->, thick, targetgreen!85!black] (ct.east) --
    node[above=3pt, font=\scriptsize\bfseries] {$2790^{2753} \bmod 3233$}
    (rec.west);

  \node[keytag=sourceblue] (pk) at (2.75, -1.1) {Public key: $(n, e) = (3233, 17)$};
  \node[keytag=structurepurple] (sk) at (8.25, -1.1) {Private key: $d = 2753$ ($e d \equiv 1 \bmod 3120$)};

  \draw[->, dashed, sourceblue!75!black] (pk.north) -- (2.75, -0.05);
  \draw[->, dashed, structurepurple!75!black] (sk.north) -- (8.25, -0.05);
\end{tikzpicture}
\end{center}

The reason it inverts -- why $(m^{e})^{d} \equiv m$ -- is Euler's theorem, and establishing that guarantee is the central goal of Chapter~19. For now, notice that everything you needed was the one operation you just learned, applied twice.

\begin{trapbox}
The numbers above are small enough to factor by trial division in well under a second: anyone who sees $n = 3233$ can recover $p = 61$ and $q = 53$, hence $\varphi(n)$, hence $d$, and read every message. That is not a flaw in RSA; it is the point. RSA's security rests entirely on $n$ being too large to factor -- real keys make $n$ a number with \emph{six hundred or more} decimal digits. A common beginner error is to admire a scheme working on toy numbers and conclude it is therefore secure. Toy numbers prove correctness, never security. Keep the two questions -- ``does it compute the right answer?'' and ``can an adversary undo it?'' -- rigidly separate in your head; conflating them is the first mistake the security mindset trains out of you.
\end{trapbox}

\section*{What it takes to be an entry-level cryptographer}

People imagine the job needs a doctorate and a gift for inventing new schemes. A few roles do. The vast majority do not. The entry-level work -- and it is real, well-paid, and in demand -- is using cryptography correctly: choosing primitives, reading specifications, reviewing designs for the failure modes that recur endlessly, integrating vetted libraries without undoing their guarantees. That work rests on a surprisingly finite base of understanding, and assembling that base is precisely what this book is for.

\begin{applicationbox}
What this looks like on the job: a protocol engineer deciding that a new service should use TLS~1.3 with X25519 key exchange and AES-256-GCM, and being able to defend each choice. A reviewer catching that a teammate reused a GCM nonce across two messages -- a one-line bug that collapses the whole scheme's integrity. An embedded developer noticing that a smartcard's signing routine branches on secret key bits and is therefore leaking through timing. None of these require inventing new mathematics. All of them require understanding the mathematics that already exists well enough to see where it is being misused.
\end{applicationbox}

\begin{recoverybox}
A fair self-test for ``ready for entry-level work,'' which doubles as a map of where you are headed. By the end you should be able to, without notes:

\begin{itemize}
  \item state Kerckhoffs's principle and explain why a public algorithm is a strength, not a weakness;
  \item compute a small RSA encryption and decryption by hand and explain why decryption inverts encryption;
  \item explain the difference between confidentiality and integrity, and why encryption alone gives you neither against an active attacker;
  \item describe a Diffie--Hellman exchange and say precisely what an eavesdropper does and does not learn from the wire;
  \item name three real-world attacks (a padding oracle, a nonce reuse, a timing leak) and sketch how each one works;
  \item read ``AES-256-GCM with a random 96-bit nonce'' and say what every part of that means and why each part is there;
  \item explain why you should reach for a vetted library and never the toy code in this book.
\end{itemize}

If those feel impossibly far off right now, good -- that is the distance the staircase covers. None of them require talent you do not have; they require the chapters ahead, worked rather than skimmed.
\end{recoverybox}

A last word on how to use the book well. \emph{Do the arithmetic.} When a worked example traces $7^{13} \bmod 23$, get a pen and follow it; when an exercise asks you to break a weakened scheme, actually break it. Cryptography punishes the reader who nods along -- the understanding that lets you spot a nonce-reuse bug in someone else's protocol is built only by having reused a nonce yourself, on small numbers, and watched the secret fall out. Run the toy code. Trace the attacks. Work through the security proofs step by step until they stop feeling like formal exercises and start reading like plain common sense.

Every chapter ends with exercises -- some to compute, some to break, some to prove. They are not optional garnish; they are where the chapter actually finishes landing. The compute problems build fluency, the break problems build the adversarial reflex that defines the field, and the prove problems build the precision that separates a guarantee from a hope. To set the tone, the preface has a few of its own. None require machinery we have not already shown; all reward a pen and ten unhurried minutes.

\section*{Exercises}

\begin{exercise}[Square-and-multiply by hand]
Compute $5^{20} \bmod 23$ using square-and-multiply, reducing mod $23$ at every step. Write $20$ in binary first, list the squarings $5^{1}, 5^{2}, 5^{4}, 5^{8}, 5^{16}$ each reduced mod $23$, then multiply together only the powers the binary digits of $20$ select. How many squarings and how many multiplications did you use, and how does that compare to the nineteen multiplications of the naive method?
\end{exercise}

\begin{exercise}[Finish the worked RSA decryption]
Using the toy key $n = 3233$, $d = 2753$ from the chapter, verify that the ciphertext $c = 2790$ decrypts to $m = 65$. Do it with square-and-multiply on $2790^{2753} \bmod 3233$: write $2753$ in binary, build the squarings $2790^{2^{k}} \bmod 3233$, and combine the selected ones. (If you have a computer handy, also confirm with \texttt{pow(2790, 2753, 3233)} -- but do the first few squaring steps by hand so the mechanism is yours.)
\end{exercise}

\begin{exercise}[Break the toy key]
The public key in the chapter is $(n, e) = (3233, 17)$. Playing the attacker, factor $n$ by trial division, recover $p$ and $q$, compute $\varphi(n)$, and reconstruct the private exponent $d$ as the inverse of $e$ modulo $\varphi(n)$. Confirm you get $d = 2753$. Then explain in one sentence why repeating this attack is hopeless when $n$ has six hundred digits rather than four.
\end{exercise}

\begin{exercise}[Scaling up modular exponentiation]
In computing $7^{13} \bmod 23$, we decomposed the exponent $13$ into binary ($1101_2$) and performed $4$ squarings and $2$ multiplications. Suppose an exponent has $2048$ bits, with roughly half of its bits set to $1$. Approximately how many squarings and multiplications would square-and-multiply require? Contrast this with the naive method of multiplying the base by itself $2^{2048}$ times.
\end{exercise}

\begin{exercise}[Eavesdroppers versus active adversaries]
Alice sends an email to Bob over open coffee-shop WiFi without encryption. Eve captures the raw radio packets. Mallory runs a rogue WiFi access point that intercepts, alters, and drops packets. For each of the following goals, state whether Eve alone can achieve it, Mallory alone can achieve it, or both: (a) reading the body of Alice's email, (b) changing the bank routing number inside the email, (c) preventing the email from reaching Bob altogether.
\end{exercise}

\begin{exercise}[Kerckhoffs in the wild]
Find a real product or protocol whose marketing claims security comes partly from a secret or proprietary algorithm. State, in the language of this preface, exactly which assumption its security is resting on, and describe the single event (a leak, a reverse-engineering, a disgruntled employee) that would collapse it. Contrast this with what a leaked \emph{key} would cost a Kerckhoffs-compliant design.
\end{exercise}

\begin{exercise}[Confidentiality is not integrity]
A naive scheme encrypts a one-byte message $m$ as $c = m \xor k$ for a secret one-byte key $k$ (a one-time pad on a single byte). Eve cannot read $m$ from $c$ without $k$. But Mallory, who controls the wire, wants to flip the message from the byte $\texttt{0x00}$ (``no'') to $\texttt{0xFF}$ (``yes'') without knowing $k$ or $m$. Show exactly what Mallory XORs into $c$ to do this, and argue that she succeeds with certainty. This is your first taste of why confidentiality alone buys you nothing against an active attacker -- and why authenticated encryption exists.
\end{exercise}

\begin{exercise}[A timing leak you can see]
Look again at the \texttt{modexp} code in this chapter. Suppose the exponent is the secret. Argue informally that two exponents with different numbers of $1$-bits will, on average, cause the function to run for measurably different lengths of time. Sketch what an attacker who can time many decryptions might learn about the secret exponent's bits. (You are not expected to mount the attack -- only to see why the door is open.)
\end{exercise}

Cryptography is, in the end, a bold mathematical claim: that two people who have never met can communicate privately across a wire the whole world is listening to, and keep their secret -- not by hiding, but by mathematics so structured that even an adversary who knows everything about the method, and has read every word of this book, still cannot break it. In Chapter~1 we start by pinning down the ground rules: who the players are, what attack models govern real systems, and what ``secure'' is actually allowed to mean.

\chapterline
